% Experimental plan: prompt composition and ablation
% Included from report.tex, Chapter 2

\section{Expérience~3~: composition et ablation de prompts}\label{sec:exp3}

L'expérience~1 utilise un prompt minimaliste de 40~mots~; l'expérience frontier
un prompt de 1\,200~mots intégrant cadrage sectoriel, narratifs par centrale,
statistiques croisées et bibliographie annotée. Or Opus obtient F1\,=\,0,64
avec le prompt court et F1\,=\,0,54 avec le prompt long~: davantage
d'instructions \emph{dégrade} l'extraction structurée. Plutôt que de
diagnostiquer ce prompt existant, nous concevons une famille de prompts
modulaires pour déterminer quels composants améliorent ou dégradent
l'extraction~:

Deux approches symétriques isolent les effets~:

\begin{description}
\item[Composition] Partir d'un prompt de base optimisé pour l'extraction CSV,
  puis ajouter un module à la fois. Mesure~: quel ajout améliore ou dégrade le
  score~?
\item[Ablation] Partir du prompt complet (base + tous les modules), puis
  retirer un module à la fois. Mesure~: quel retrait améliore ou dégrade le
  score~?
\end{description}

\noindent La comparaison des deux bras révèle les \emph{effets d'interaction}~:
un module peut aider seul mais nuire en combinaison, ou inversement.

\subsection{Prompt de base}\label{sec:prompt-base}

Le prompt de base a été conçu par itérations successives avec les critères
suivants~: (a)~demander un CSV structuré avec colonnes nommées correspondant
au parseur d'évaluation, (b)~spécifier un seuil d'inclusion (30~MWe),
(c)~couvrir tout le cycle de vie des projets (proposition à démantèlement),
(d)~inclure les centrales captives et de cogénération industrielle,
(e)~exclure explicitement la biomasse et les déchets (suivis séparément),
(f)~utiliser un vocabulaire de statut ouvert reflétant la multiplicité des
autorités décisionnaires (investisseurs, provinces, ministère, accords
internationaux), (g)~ne pas inclure de persona (la littérature sur le
\emph{role prompting} suggère un effet négatif sur l'extraction structurée),
(h)~ne pas révéler le nombre attendu de résultats.

\medskip
\noindent La terminologie est alignée sur l'Open Energy Ontology (OEO)
avec trois divergences délibérées~: inclusion des centrales hors réseau
(l'OEO exige la connexion au réseau), distinction gaz domestique / GNL
importé (l'OEO ne distingue pas la provenance), et regroupement «~thermique~»
par usage commun plutôt que par classe OEO formelle.

\label{prompt:base}%
\begin{Prompt}
Produce a comprehensive CSV table listing ALL thermal power plants in
Vietnam with capacity of 30 MWe or above, at every stage from early
proposal through decommissioning. Include grid-connected utility plants,
captive (self-generation) plants, and industrial cogeneration facilities.

Columns:
- Name_VI: plant name in Vietnamese
- Name_EN: plant name in English
- Fuel: the primary fuel — e.g. Coal, Domestic gas, Imported LNG, Oil.
  Exclude biomass and waste-to-energy plants (tracked separately).
- Technology: Subcritical / Supercritical / USC / CCGT / OCGT / CFB
- Status: current stage of the project, e.g. Operational, Under
  construction, Proposed, Permitted, Financed, Suspended, Cancelled,
  Decommissioned. Use the most specific term that applies. Projects may
  have been proposed by investors, provincial authorities, international
  partners, or national planners — capture all regardless of which
  authority originated them.
- COD: actual or expected commercial operation date
- Province: location
- Capacity_MWe: current or as-built net generation capacity in MWe.
  If the plant was proposed at a different capacity, note the difference.
- Owner: key investors or controlling entities — e.g. "EVN/GENCO 1",
  "Marubeni + KEPCO (BOT)", "Formosa Plastics". Skip single-purpose
  project vehicles.
- Note: additional context — e.g. original proposed capacity if different
  from as-built, unit breakdowns, phase information, ownership changes.

Be exhaustive. Cover every thermal power plant you know of, whether it
appeared in a national power development plan, regulatory filing,
investment proposal, international agreement, news report, or any other
source. Include cancelled projects, suspended proposals, and plants that
were planned but never built. Missing plants reduces the value of this
inventory.

Format the response as CSV with a header row.
\end{Prompt}

\subsection{Modules}\label{sec:modules}

Six modules sont testés. Chacun est un paragraphe ajouté à la fin du prompt
de base.

\paragraph{Persona} Ajoute une consigne de rôle en tête du prompt.
\label{prompt:mod-persona}%
\begin{Prompt}
You are a senior energy analyst.
\end{Prompt}

\paragraph{Cadrage sectoriel} Demande un paragraphe introductif avant le CSV.
\label{prompt:mod-overview}%
\begin{Prompt}
Before the CSV table, provide a sector overview covering: Vietnam's
electricity generation mix (installed capacity, generation shares by
fuel), key policy documents and their thermal targets, fuel supply
landscape (domestic coal, domestic gas, LNG import projects), and key
institutional actors (state utilities, gas companies, BOT investors,
provincial authorities).
\end{Prompt}

\paragraph{Narratifs par centrale} Demande un commentaire sourcé après le CSV.
\label{prompt:mod-narratives}%
\begin{Prompt}
After the CSV table, provide a sourced paragraph for each plant complex
covering: development history and timeline, notable issues (delays,
financing changes, ownership transfers, technology changes), and status
confidence (HIGH = government decision or company report, MEDIUM =
multiple news sources, LOW = older plans only, status uncertain).
\end{Prompt}

\paragraph{Bibliographie} Demande une liste annotée de sources.
\label{prompt:mod-bibliography}%
\begin{Prompt}
After completing the inventory, list every source you relied on as an
annotated bibliography organized by category (government decisions,
power development plans, utility reports, company filings, news). For
each source: full citation, URL if known ("URL not verified" if
unsure), and one sentence on what information was drawn from it.
\end{Prompt}

\paragraph{Tableaux statistiques} Demande des tableaux croisés récapitulatifs.
\label{prompt:mod-statistics}%
\begin{Prompt}
After the inventory, produce cross-tabulation summary tables:
a) Capacity by fuel x status (total MWe per cell, with row and column
   totals)
b) Top 15 provinces by total thermal capacity, broken down by fuel
c) Timeline of capacity additions by period and fuel type
\end{Prompt}

\paragraph{Sourçage} Demande des colonnes de citation par ligne et une liste
de références.
\label{prompt:mod-sourcing}%
\begin{Prompt}
For each row, include Source_1 and Source_2 columns citing the primary
source for that plant's data. After the table, list all cited sources
with full citations. Do not fabricate URLs — write "URL not verified"
if unsure of the exact address.
\end{Prompt}

\subsection{Plan de mesures}\label{sec:plan-mesures}

Le design expérimental comporte deux bras symétriques et trois points d'ancrage.

\paragraph{Points d'ancrage}

\begin{description}
\item[Census] Le Prompt~1 original (40~mots), point de référence bas.
\item[Base] Le prompt de base (section~\ref{sec:prompt-base}), point de
  référence intermédiaire.
\item[Complet] Le prompt de base + les 6 modules, point de référence haut.
\item[Frontier] Le prompt frontier existant (1\,200~mots), contrôle externe.
\end{description}

\paragraph{Bras composition} À partir du prompt de base, ajouter un seul module
à la fois (6~variantes).

\paragraph{Bras ablation} À partir du prompt complet, retirer un seul module
à la fois (6~variantes).

\medskip
\noindent Soit 16~prompts au total~:

\begin{center}
\small
\begin{tabular}{rllc}
\toprule
\textbf{\#} & \textbf{Prompt} & \textbf{Bras} & \textbf{Modules} \\
\midrule
 0 & Census (Prompt~1 existant) & Ancrage  & -- \\
 1 & Base                       & Ancrage  & $\emptyset$ \\
 2 & Base + persona             & Composition & \{P\} \\
 3 & Base + cadrage             & Composition & \{O\} \\
 4 & Base + narratifs           & Composition & \{N\} \\
 5 & Base + bibliographie       & Composition & \{B\} \\
 6 & Base + statistiques        & Composition & \{T\} \\
 7 & Base + sourçage            & Composition & \{S\} \\
 8 & Complet                    & Ancrage  & \{P,O,N,B,T,S\} \\
 9 & Complet $-$ persona        & Ablation & \{O,N,B,T,S\} \\
10 & Complet $-$ cadrage        & Ablation & \{P,N,B,T,S\} \\
11 & Complet $-$ narratifs      & Ablation & \{P,O,B,T,S\} \\
12 & Complet $-$ bibliographie  & Ablation & \{P,O,N,T,S\} \\
13 & Complet $-$ statistiques   & Ablation & \{P,O,N,B,S\} \\
14 & Complet $-$ sourçage       & Ablation & \{P,O,N,B,T\} \\
15 & Frontier (existant)        & Contrôle & -- \\
\bottomrule
\end{tabular}
\end{center}

\subsection{Protocole}\label{sec:protocole-ablation}

\paragraph{Phase de sélection} Les prompts~1 (base) et~8 (complet) sont
exécutés sur l'ensemble des modèles du census (25+ modèles cloud) avec 2~runs
chacun. On retient les $k$~modèles dont le $|\Delta\text{F1}|$ moyen entre
les deux conditions dépasse 0,05 (5~points). Ce seuil élimine les modèles
insensibles au prompt, pour lesquels l'ablation ne peut rien révéler.

\paragraph{Phase de mesure} Les 16~prompts sont exécutés sur les $k$~modèles
retenus, avec 2~répétitions chacun ($16 \times k \times 2$ runs). Température
réglée selon les recommandations du fournisseur pour l'extraction structurée
(0,0--0,3 selon les modèles~; pour DeepSeek, la température API 1,0
correspond à une température interne de 0,3 via un mécanisme de mapping
documenté). Le nombre total de tokens générés est mesuré par condition pour
contrôler l'effet de budget de sortie (un prompt demandant des narratifs
produit mécaniquement plus de tokens, ce qui peut réduire la place disponible
pour l'inventaire). Chaque run produit un CSV évalué contre l'inventaire de
référence (163~centrales, matching par nom exact et approximatif).

\paragraph{Analyse}
\begin{itemize}
\item \textbf{Effets principaux}~: pour chaque module $m$, comparer
  F1(base) \emph{vs} F1(base+$m$) (composition) et F1(complet) \emph{vs}
  F1(complet$-m$) (ablation). Les intervalles de confiance bootstrap à 95\,\%
  sur le $\Delta$F1 déterminent la significativité.
\item \textbf{Effets d'interaction}~: si un module aide en composition mais
  nuit en ablation (ou inversement), il interagit avec les autres modules.
  L'asymétrie composition/ablation quantifie cet effet.
\item \textbf{Coût de la confiance}~: le module sourçage est non négociable
  pour le pipeline de production (section~\ref{chap:conclusion}). Son
  $\Delta$F1 mesure le prix, en points d'extraction, de la traçabilité.
\end{itemize}

\paragraph{Budget prévisionnel (pilote)} Sélection~: $2 \times 25 \times 2
= 100$~runs. Mesure avec $k=5$ modèles~: $16 \times 5 \times 2 = 160$~runs.
Total~: 260~runs $\approx$ 5--15\,\$. Les modèles open-weight réduisent ce
budget d'un facteur~10. Pour une publication, augmenter à 5~répétitions,
élargir le panel de modèles, et pré-enregistrer les hypothèses directionnelles
par module (section~\ref{sec:hypotheses-ablation}).

\paragraph{Limites du design}
\begin{itemize}
\item \textbf{Requête directe uniquement.} L'expérience teste les prompts
  en requête single-shot, sans RAG ni décomposition. L'interaction entre
  le design du prompt et le contexte RAG constitue une expérience séparée~:
  un module qui nuit en single-shot (cadrage sectoriel gaspillant des tokens)
  pourrait aider en RAG (amorçage de la récupération documentaire).
\item \textbf{Plan fractionnaire (résolution~III).} Le design teste les
  effets principaux de 6~modules (16~prompts) mais pas les $\binom{6}{2}=15$
  interactions par paires (qui nécessiteraient un plan factoriel $2^6=64$
  prompts). L'asymétrie composition/ablation détecte la \emph{présence}
  d'interactions mais ne les attribue pas à un couple spécifique.
\end{itemize}

\subsection{Hypothèses pré-enregistrées}\label{sec:hypotheses-ablation}

Les hypothèses ci-dessous sont formulées \emph{avant} l'analyse des données
d'ablation. Elles s'appuient sur trois mécanismes théoriques issus de la
littérature sur l'extraction structurée par LLM~:

\begin{enumerate}
\item \textbf{Compétition de budget de sortie.} Tout module demandant du texte
  non tabulaire (narratifs, bibliographie, cadrage) consomme des tokens qui
  auraient pu produire des lignes CSV. Ce mécanisme dégrade le rappel sans
  affecter directement la précision.
\item \textbf{Amorçage par contexte (\emph{priming}).} Un module qui active
  des connaissances pertinentes (noms de centrales, vocabulaire sectoriel) peut
  améliorer le rappel en rendant ces entrées plus saillantes dans la génération.
\item \textbf{Surcharge d'instructions (\emph{instruction overload}).} Un
  prompt trop long dilue les consignes d'extraction structurée, ce qui peut
  dégrader la précision (erreurs de format, colonnes manquantes).
\end{enumerate}

\paragraph{Hypothèses par module}

\begin{center}
\small
\begin{tabular}{l c c p{0.52\textwidth}}
\toprule
\textbf{Module} & \textbf{Rappel} & \textbf{Précision} & \textbf{Mécanisme attendu} \\
\midrule
Persona       & neutre  & neutre
  & Un rôle générique (\emph{senior energy analyst}) n'ajoute ni
    connaissances ni contraintes. La littérature sur le \emph{role
    prompting} montre un effet nul ou négatif sur l'extraction
    structurée. \\
\addlinespace
Cadrage       & $+$     & neutre
  & Le résumé sectoriel amorce le vocabulaire du domaine (noms de
    centrales, provinces, acteurs) et rend des entrées marginales
    plus saillantes. Coût de budget modéré ($\sim$200~tokens). \\
\addlinespace
Narratifs     & $-$     & neutre
  & Demander un paragraphe par complexe industriel consomme
    ${>}\,2\,000$~tokens de sortie. Mécanisme dominant~: compétition
    de budget. Effet de rappel négatif attendu le plus fort. \\
\addlinespace
Bibliographie & $-$     & neutre
  & La liste de sources annotées consomme ${\sim}\,1\,000$~tokens de
    sortie. Même mécanisme de compétition de budget que les
    narratifs, mais moindre car le texte est plus compact. \\
\addlinespace
Statistiques  & $-$     & neutre
  & Les tableaux croisés consomment ${\sim}\,500$~tokens de sortie.
    Compétition de budget modérée. Effet attendu faible. \\
\addlinespace
Sourçage      & $-$     & $+$
  & Ajouter Source\_1 et Source\_2 par ligne force le modèle à
    justifier chaque entrée, ce qui filtre les hallucinations
    (précision~$+$). En contrepartie, le format élargi (12~colonnes
    au lieu de~10) réduit le nombre de lignes produites (rappel~$-$). \\
\bottomrule
\end{tabular}
\end{center}

\paragraph{Hypothèses composites}

\begin{enumerate}
\item \textbf{Base \emph{vs} Census.} Le prompt de base est attendu supérieur
  au census (prompt~1 de 40~mots) en rappel et en précision. Le prompt de base
  spécifie le format CSV, les colonnes exactes, le seuil de 30~MWe et les
  critères d'inclusion/exclusion~; le census ne fournit qu'une consigne
  minimale. $H_1 : \text{F1}_{\text{base}} > \text{F1}_{\text{census}}$.

\item \textbf{Complet \emph{vs} Base.} Le prompt complet (6~modules) est
  attendu inférieur au prompt de base en rappel, en raison de la compétition
  de budget cumulée (narratifs + bibliographie + statistiques + sourçage).
  L'effet d'amorçage du cadrage ne compense pas le coût combiné.
  $H_2 : \text{Rappel}_{\text{complet}} < \text{Rappel}_{\text{base}}$.

\item \textbf{Complet \emph{vs} Frontier.} Le prompt complet et le prompt
  frontier sont de longueur comparable ($\sim$1\,200~mots). Le prompt complet
  est attendu supérieur car sa base a été optimisée pour l'extraction CSV
  structurée, alors que le frontier mêle instructions et contenu éditorial.
  $H_3 : \text{F1}_{\text{complet}} > \text{F1}_{\text{frontier}}$.

\item \textbf{Module à effet individuel le plus fort.} Les narratifs sont
  attendus comme le module ayant le plus grand $|\Delta\text{F1}|$ individuel,
  car ils imposent la plus forte compétition de budget de sortie
  ($>$2\,000~tokens de texte libre demandé). Le sourçage est le second
  candidat, car il modifie à la fois le format de sortie et le processus de
  génération.
\end{enumerate}

\paragraph{Analyse empirique Base \emph{vs} Census (H$_1$)}

L'intersection des deux bras couvre désormais neuf modèles (DeepSeek-V3.2,
Mistral-Small-2603, Gemini-2.5-Flash-Lite, Gemini-3-Flash-Preview,
Kimi-K2.5, Minimax-M2.7, Qwen3.5-122b-A10b, Qwen3.5-35b-A3b et
Seed-2.0-Lite), après avoir étendu la campagne Phase~1 à sept modèles
\emph{cloud} bon marché supplémentaires déjà présents dans le panel
\emph{census}. L'échantillon reste modeste mais autorise une décomposition
quantitative.

% Auto-generated by aedist.tabulate_base_vs_census — do not edit
\begin{table}[htbp]
\centering
\caption{Prompt \emph{base} (300~mots, structuré) \emph{vs} prompt \emph{census} (40~mots, minimal) — moyennes par modèle sur l'intersection des deux bras.}\label{tab:base-vs-census}
\begin{tabular}{@{}lrrrrrrrr@{}}
\toprule
 & \multicolumn{3}{c}{Census} & \multicolumn{3}{c}{Base} & & \\
\cmidrule(lr){2-4}\cmidrule(lr){5-7}
Modèle & P & R & F1 & P & R & F1 & $\Delta$F1 (pp) & $\Delta$tok$_\text{in}$ \\
\midrule
Seed-2.0-Lite & 100.0\% & 7.8\% & 14.3\% & 100.0\% & 23.9\% & 38.6\% & +24.3 & +379 \\
DeepSeek-V3.2 & 100.0\% & 17.8\% & 30.1\% & 100.0\% & 42.3\% & 59.4\% & +29.3 & +362 \\
Gemini-2.5-Flash-Lite & 97.5\% & 53.0\% & 62.3\% & 35.6\% & 100.0\% & 51.7\% & $-$10.6 & +378 \\
Gemini-3-Flash-Preview & 100.0\% & 33.7\% & 50.4\% & 100.0\% & 49.1\% & 65.8\% & +15.4 & +378 \\
Minimax-M2.7 & 66.7\% & 17.0\% & 27.0\% & 100.0\% & 45.1\% & 61.2\% & +34.3 & +354 \\
Mistral-Small-2603 & 100.0\% & 30.3\% & 46.3\% & 100.0\% & 33.7\% & 49.5\% & +3.2 & +365 \\
Kimi-K2.5 & 100.0\% & 36.8\% & 53.4\% & 100.0\% & 60.7\% & 75.6\% & +22.1 & +357 \\
Qwen3.5-122b-A10b & 100.0\% & 24.5\% & 39.1\% & 100.0\% & 44.2\% & 61.3\% & +22.2 & +373 \\
Qwen3.5-35b-A3b & 100.0\% & 20.9\% & 34.3\% & 32.8\% & 100.0\% & 49.4\% & +15.1 & +373 \\
\midrule
\multicolumn{7}{r}{\textit{Macro-moyenne $\Delta$F1 (IC 95\,\% bootstrap)}} & +17.3 [+8.4, +25.1] & \\
\bottomrule
\multicolumn{9}{l}{\footnotesize\textit{n{=}9 modèles dans l'intersection des deux bras~; underpowered pour inférence générale, valeur illustrative. Tokens d'entrée = moyenne par run.}} \\
\end{tabular}
\end{table}


\begin{figure}[htbp]
\centering
\includegraphics[width=0.75\linewidth]{inputs/generated/fig_base_vs_census.pdf}
\caption{F1 par modèle sur le prompt \emph{base} \emph{vs} le prompt
\emph{census}. La diagonale $y=x$ matérialise l'absence d'effet~; la taille
des points est proportionnelle à $\log(\Delta\text{tok}_\text{in} + 1)$.}
\label{fig:base-vs-census}
\end{figure}

La macro-moyenne mesurée est $\Delta\text{F1} = +17{,}3$~pp avec un
intervalle de confiance bootstrap 95\,\% de $[+8{,}4\,;\,+25{,}1]$~pp
(n${=}9$, graine fixée). L'IC exclut zéro dans la direction prédite~:
\textbf{H$_1$ est corroborée}. Huit modèles sur neuf affichent un gain
positif (de $+3{,}2$~pp pour Mistral-Small-2603 à $+34{,}3$~pp pour
Minimax-M2.7)~; le seul cas négatif est Gemini-2.5-Flash-Lite à
$-10{,}6$~pp, où le bras \emph{base} a saturé la fenêtre de sortie
(\emph{finish=length}) en générant 476 faux positifs et faisant chuter la
précision à 35{,}6\,\%. La décomposition précision/rappel est éclairante~:
pour la majorité des modèles (DeepSeek, Mistral, Gemini-3-Flash, Kimi,
Qwen3.5-122b, Seed) la précision reste à 100\,\% sur les deux bras et
\emph{l'intégralité du gain en F1 provient d'un meilleur rappel} — le
format colonné et l'instruction d'exhaustivité agissent comme
amplificateurs de rappel plutôt que comme filtres d'hallucinations.
Cependant, sur deux modèles (Gemini-2.5-Flash-Lite, Qwen3.5-35b-A3b) le
prompt \emph{base} déclenche une dégradation de la précision~: le modèle
tente de remplir des colonnes qu'il devine, au lieu d'omettre les valeurs
manquantes. L'effet est donc fortement \emph{model-dependent} et
non-monotone~: les modèles déjà calibrés tirent un bénéfice de rappel sans
coût de précision, tandis que certains petits modèles ou modèles
expérimentaux peuvent payer un coût de précision en suivant les
instructions plus prescriptives. Nous différons à un travail ultérieur la
vérification de robustesse au seuil 30~MWe (retirer explicitement le seuil
du prompt \emph{base} et re-tester) qui isolerait la contribution
spécifique de cette contrainte au sein de l'écart observé.

\paragraph{Prompt census \emph{vs} prompt base~: deux objets différents}
L'écart mesuré entre census et base ne relève pas du bruit~: il reflète
une distinction fonctionnelle. Le prompt \emph{census} (40~mots, sans
seuil, sans vocabulaire de statut) mesure la capacité brute d'extraction
du modèle — c'est l'analogue d'un test de QI pour l'extraction
structurée. Le prompt \emph{base} (280~mots, 9~colonnes, seuil 30~MWe,
règles d'inclusion/exclusion) mesure l'utilisabilité dans un pipeline de
production. L'écart de F1 entre les deux est donc le \emph{prix de la
mise en production}~: les tokens d'instruction supplémentaires
($\Delta\text{tok}_\text{in} \approx +370$ par appel) achètent un
rappel accru sans dégradation de précision chez la majorité des modèles.
Ce coût est modeste ($< 0{,}001$~USD par appel aux tarifs actuels), mais
se multiplie par le nombre de campagnes et de répétitions dans un
déploiement à l'échelle (60~campagnes $\times$ 3~répétitions = 180
appels dans le scénario Auto-PyPSA ASEAN). Le prompt \emph{census}
reste pertinent comme référence de capacité minimale~; le prompt
\emph{base} est l'ancrage opérationnel pour les comparaisons
d'ablation.

\paragraph{Seuils de falsification} Une hypothèse directionnelle est réfutée
si l'intervalle de confiance bootstrap à 95\,\% sur $\Delta$F1 exclut zéro
dans la direction opposée à la prédiction. Une hypothèse «~neutre~» est
réfutée si $|\Delta\text{F1}| > 0{,}03$ (3~points) avec un IC à 95\,\%
excluant zéro.

\subsection{Régimes d'information et méthodes du benchmark}\label{sec:regime-linkage}

Le benchmark présenté aux chapitres~\ref{chap:evaluation} et~\ref{chap:rag}
évalue cinq méthodes d'extraction sur un panel de modèles~:

\begin{enumerate}
\item \textbf{Requête directe} (section~\ref{sec:exp1})~: un seul prompt,
  sans document externe ni recherche.
\item \textbf{Multi-tour} (section~\ref{sec:exp2})~: conversation avec
  relances successives.
\item \textbf{RAG wholesale} (section~\ref{sec:rag})~: injection de
  l'intégralité d'un corpus de 18~documents dans le contexte.
\item \textbf{Décomposition} (section~\ref{sec:decomposition})~: découpage
  de la requête par type de combustible, puis fusion.
\item \textbf{Recherche web}~: le modèle accède à Internet via un outil
  de recherche (Tavily ou plugin natif).
\end{enumerate}

L'expérience d'ablation (section~\ref{sec:plan-mesures}) croise 16~variantes
de prompt avec trois \emph{régimes d'information}, c'est-à-dire trois modes
d'accès aux connaissances~:

\begin{description}
\item[Paramétrique] Le modèle répond uniquement à partir de ses poids
  pré-entraînés. Ce régime correspond à la \textbf{requête directe}.
\item[RAG] Le modèle reçoit le corpus de 18~documents en contexte.
  Ce régime correspond au \textbf{RAG wholesale}.
\item[Recherche web] Le modèle dispose d'un outil de recherche Internet.
  Ce régime correspond à la \textbf{recherche web}.
\end{description}

\noindent Les deux méthodes restantes ne constituent pas des régimes
d'information et sont traitées par des expériences dédiées~:

\paragraph{Multi-tour $\to$ vérification.}
L'approche conversationnelle du benchmark utilise des relances génériques
(«~Avez-vous oublié des centrales~?~»). Or l'expérience de vérification
multi-agents (section~\ref{sec:exp5}) subsume cette approche en remplaçant
les relances aveugles par des vérifications ciblées --- chaque tour a un
objectif spécifique (contrôler la provenance, détecter les hallucinations,
croiser avec le référentiel). Le multi-tour générique est donc un cas
particulier, moins motivé, du protocole de vérification.

\paragraph{Décomposition $\to$ exclue (axe orthogonal).}
La décomposition par combustible (section~\ref{sec:decomposition}) modifie
la \emph{granularité de la requête}, non le régime d'information~: on peut
décomposer en paramétrique, en RAG ou en recherche web. C'est un axe
orthogonal. De plus, la décomposition introduit des hallucinations non
thermiques (centrales hydroélectriques ou éoliennes ajoutées dans la
catégorie «~autres combustibles~»), ce qui en fait un facteur confondant
plutôt qu'un régime à tester. L'interaction
décomposition~$\times$~régime constitue un travail futur.

\medskip
\noindent En résumé, l'ablation couvre trois des cinq méthodes du benchmark
en les reformulant comme des régimes d'information. Cette reformulation
permet de mesurer l'effet du prompt \emph{à régime fixe}~: un module qui
aide en paramétrique pourrait nuire en RAG (où le contexte documentaire
rend le cadrage sectoriel redondant), et inversement. Le croisement
prompt~$\times$~régime est le design central de cette expérience.

\subsection{Lien avec la vérification multi-agents}\label{sec:lien-verification}

Cette expérience mesure l'effet du prompt sur l'extraction en un seul tour. La
vérification --- un agent indépendant contrôle la sortie d'un autre ---
constitue une expérience séparée. La littérature sur la self-consistency
et les protocoles de débat montre que la vérification par
agents multiples et indépendants est plus efficace que l'auto-révision par le
même modèle. Le module sourçage (colonnes Source\_1, Source\_2) est conçu pour
alimenter cette phase~: chaque ligne du CSV porte ses propres références,
permettant à un agent vérificateur de contrôler la provenance sans accéder au
prompt original.

\subsection{Résultats de l'ablation --- régime RAG}\label{sec:results-ablation}

La Phase~2 teste 16~variantes de prompt sur le régime RAG avec deux modèles
frontière (DeepSeek~V3.2 et Kimi~K2~Thinking), choisis pour leurs effets
opposés en Phase~1 ($-27{,}4$\,pp et $+25{,}8$\,pp respectivement avec le
prompt composite). Chaque variante est répétée deux fois (64~appels au total,
coût~: 3{,}77\,\$).

La figure~\ref{fig:ablation-strip} montre les résultats par variante sous forme
de strip~plot~: chaque point bleu est une centrale correctement identifiée,
chaque point orange une hallucination. DeepSeek~V3.2 refuse (tool\_calls) sur
16/32~variantes, rendant l'analyse $\Delta$F1 incomplète pour ce modèle.

\begin{figure}[htbp]
\centering
\includegraphics[width=0.85\linewidth]{inputs/generated/fig_ablation_strip.pdf}
\caption{Strip plot de l'ablation Phase~2 en régime RAG. Bleu~: centrales
correctement identifiées (TP). Orange~: hallucinations (FP). Croix grises~:
refus du modèle. Ligne verte~: inventaire complet (163~centrales).}
\label{fig:ablation-strip}
\end{figure}

La figure~\ref{fig:ablation-heatmap} résume les effets par module sous forme de
matrice. Pour chaque module, deux $\Delta$F1 sont calculés~: l'effet de
composition (base + module vs base seul) et l'effet d'ablation (composite $-$
module vs composite complet).

\begin{figure}[htbp]
\centering
\includegraphics[width=0.7\linewidth]{inputs/generated/fig_ablation_heatmap.pdf}
\caption{Effet de chaque module sur le F1 en régime RAG. Rouge~: le module
dégrade les performances. Bleu~: le module les améliore. Tirets~: données
insuffisantes (refus du modèle).}
\label{fig:ablation-heatmap}
\end{figure}

Pour Kimi~K2~Thinking, le module \emph{sourçage} est le seul à améliorer
le F1 en composition ($+3{,}6$\,pp) et son retrait dégrade le composite
($-7{,}1$\,pp). Le module \emph{statistiques} a l'effet inverse~:
$-30{,}4$\,pp en composition, $+11{,}5$\,pp en ablation. L'interprétation
est que le module statistiques (format tabulaire imposé) interfère avec le
contexte RAG déjà structuré, tandis que le sourçage (demande de références)
aide le modèle à ancrer ses réponses dans les documents fournis.

\paragraph{Limitation~: température non contrôlée}
Les 64~appels de la Phase~2 ont été exécutés avec la température par défaut du
fournisseur (\texttt{temperature=null}, typiquement 0{,}7--1{,}0 selon le
modèle) plutôt qu'à $t=0$. Ce paramétrage introduit un bruit stochastique
supplémentaire entre les répétitions, mais ne biaise pas la comparaison
base/composite puisque les deux bras partagent la même configuration de
température. Les intervalles de confiance rapportés sont donc potentiellement
plus larges qu'ils ne le seraient à $t=0$, ce qui rend nos tests plus
conservateurs. Le code a depuis été corrigé pour imposer $t=0{,}0$ par défaut
(PR\,\#244)~; une reprise ciblée des modules ambigus pourra être effectuée si
les relecteurs le demandent.
